% TEMPLATE-MILC-v3.tex - plantilla maestra UNICA de las guias MILC v3.
%
% La usan las tres familias de guias, cada una con su driver:
%   · Grados 6-11  -> scripts/build-guias-g11.py
%   · Bebras       -> scripts/build-guias-bebras.py
%   · Semillero    -> scripts/build-guias-semillero.py
%
% Antes habia tres copias casi identicas (~400 lineas iguales, 6 distintas):
% cada mejora habia que aplicarla tres veces, y en la practica no se hacia
% (el motor de recursos vivia solo en dos de ellas). Lo que varia entre
% familias esta parametrizado en 6 placeholders que cada driver llena
% (se nombran aqui SIN su sintaxis real: el builder reemplaza por texto plano,
% tambien dentro de los comentarios, y un valor multilinea rompe el preambulo):
%
%   HEADER_IZQ        encabezado izquierdo de pagina
%   HEADER_DER        encabezado derecho de pagina
%   PORTADA_ETIQUETA  rotulo sobre el numero grande (GRADO/MOMENTO/MODULO)
%   PORTADA_SERIE     linea de serie de la portada
%   PORTADA_ID        identificador de la guia en la portada
%   PORTADA_CIRCULO   el circulito de grado (°), o vacio si no aplica
%
% El resto de claves las documenta content/guias/_SCHEMA.md. El script valida
% que no queden placeholders sin reemplazar antes de escribir el archivo final.

\documentclass[11pt,letterpaper]{article}
\usepackage[margin=1.75cm]{geometry}
\usepackage[table]{xcolor}
\usepackage{fontspec}
\usepackage[spanish]{babel}
\usepackage{tikz}
% Librerías TikZ para los diagramas de las guías (bloque `recursos:`):
%  · babel  → babel en español vuelve activos caracteres como «>», y TikZ los usa
%             en sus opciones (>=stealth, ->). Sin esta librería, cualquier
%             diagrama con flechas revienta con «\language@active@arg> has an
%             extra }». Debe ir ANTES de que se dibuje nada.
%  · positioning        → sintaxis «below left=2pt and 3pt».
%  · decorations.pathreplacing → llaves ({brace}) para señalar rangos.
\usetikzlibrary{babel,positioning,decorations.pathreplacing}
\usepackage{graphicx}
% Circuitos: el trabajo de electrónica se hace hoy con micro:bit y sus sensores
% integrados (MakeCode, sin cableado), así que no se necesitan esquemáticos.
% Si algún día entran montajes con protoboard, basta descomentar esta línea:
% los diagramas .tex/.tikz declarados en `recursos:` ya se incrustan solos.
% \usepackage{circuitikz}
\usepackage[most]{tcolorbox}
\usepackage{tabularx}
\usepackage{array}
\usepackage{fancyhdr}
\usepackage{titlesec}
\usepackage{ragged2e}
\usepackage{enumitem}
\usepackage{microtype}

% Helvetica Neue no trae la flecha U+2192, asi que cada «→» escrito literalmente
% en el YAML se compilaba a NADA, en silencio (462 flechas en 61 guias: rutas del
% tipo «paleta → tipografia → lockup» salian rotas). Mapeamos el caracter a su
% version matematica, que si tiene glifo. Asi el autor puede seguir escribiendo
% «→» comodamente en el YAML.
\usepackage{newunicodechar}
\newunicodechar{→}{\ensuremath{\rightarrow}}

\setmainfont{Helvetica Neue}
\setsansfont{Helvetica Neue}
\setlength{\parindent}{0pt}
\setlength{\parskip}{6pt}
\hyphenpenalty=200
\exhyphenpenalty=200
\pretolerance=400
\tolerance=2000
\setlength{\emergencystretch}{1em}
\renewcommand{\arraystretch}{1.25}
\setlist{leftmargin=5mm,itemsep=1mm,topsep=1mm}
\newcolumntype{Y}{>{\RaggedRight\arraybackslash}X}

\definecolor{milcMagenta}{HTML}{E6007E}
\definecolor{milcVino}{HTML}{5A0038}
\definecolor{milcTurquesa}{HTML}{0093A5}
\definecolor{milcVerde}{HTML}{58A923}
\definecolor{milcMostaza}{HTML}{E5B400}
\definecolor{milcOcre}{HTML}{B86600}
\definecolor{milcNegro}{HTML}{241816}
\definecolor{milcGris}{HTML}{F7F7F4}
\definecolor{milcLinea}{HTML}{E8E2DD}
\definecolor{milcGradePrimary}{HTML}{84CC16}
\definecolor{milcGradeDark}{HTML}{4D7C0F}
\definecolor{milcGradeSoft}{HTML}{ECFCCB}
\definecolor{milcGradeAccent}{HTML}{CFE7FF}
\definecolor{milcGradeText}{HTML}{FFFFFF}
\color{milcNegro}

\pagestyle{fancy}
\fancyhf{}
\lhead{\small Tecnología e Informática · Bebras}
\rhead{\small Momento 1 · MILC v3}
\cfoot{\small\thepage}
\renewcommand{\headrulewidth}{0pt}

\newtcolorbox{titlebox}[1]{
  enhanced,colback=#1,colframe=#1,arc=7mm,boxrule=0pt,
  left=7mm,right=7mm,top=3mm,bottom=3mm,
  before skip=8pt,after skip=8pt,
  fontupper=\sffamily\bfseries\Large\color{white}
}

\newtcolorbox{softbox}[3]{
  enhanced,colback=#2,colframe=#1,borderline west={4pt}{0pt}{#1},
  arc=2mm,boxrule=.4pt,left=4mm,right=4mm,top=3mm,bottom=3mm,
  before skip=6pt,after skip=8pt,title={#3},coltitle=white,
  colbacktitle=#1,fonttitle=\sffamily\bfseries,fontupper=\RaggedRight,
  attach boxed title to top left={xshift=3mm,yshift=-2mm},
  boxed title style={arc=2mm, boxrule=0pt}
}

\newtcolorbox{drawbox}[2]{
  enhanced,colback=white,colframe=#1,arc=4mm,boxrule=1pt,
  left=5mm,right=5mm,top=4mm,bottom=4mm,before skip=8pt,after skip=8pt,
  title={#2},coltitle=white,colbacktitle=#1,fonttitle=\sffamily\bfseries,
  fontupper=\RaggedRight,
  attach boxed title to top left={xshift=4mm,yshift=-2mm},
  boxed title style={arc=3mm, boxrule=0pt}
}

\newtcolorbox{infoband}[1]{
  enhanced,colback=#1!8,colframe=#1!50,arc=2mm,boxrule=.5pt,
  left=4mm,right=4mm,top=2mm,bottom=2mm,before skip=4pt,after skip=4pt,
  fontupper=\small\RaggedRight
}

% Puente narrativo entre secciones (contrato editorial Conéctate).
% Da continuidad: "Acabas de X. Ahora vas a Y porque Z."
% Visualmente: banda izquierda en color del grado, texto en itálica pequeña.
\newtcolorbox{puentebox}{
  enhanced,colback=milcGradeSoft,colframe=milcGradeDark,
  borderline west={3pt}{0pt}{milcGradeDark},
  arc=0mm,boxrule=0pt,left=5mm,right=4mm,top=2.5mm,bottom=2.5mm,
  before skip=4pt,after skip=8pt,
  fontupper=\itshape\small\color{milcNegro}\RaggedRight
}
% La flecha va en modo matematico a proposito: Helvetica Neue NO tiene el glifo
% U+2192, asi que \textrightarrow se compilaba a NADA («Missing character: There
% is no → in font Helvetica Neue») y el puente salia sin flecha. En modo
% matematico la flecha viene de la fuente matematica y si se dibuja.
\newcommand{\puente}[1]{%
  \begin{puentebox}%
    \textcolor{milcGradeDark}{$\rightarrow$}\hspace{2mm}#1%
  \end{puentebox}%
}

\newcommand{\linead}[1]{\noindent\color{milcLinea}\rule{#1}{0.5pt}\color{milcNegro}}
\newcommand{\respuesta}[1]{\par\vspace{2mm}\foreach \n in {1,...,#1}{\noindent\color{milcLinea}\rule{\linewidth}{0.5pt}\par\vspace{5mm}}\color{milcNegro}}
\newcommand{\checkbox}{\textcolor{milcGradeDark}{\fbox{\phantom{x}}}\hspace{5pt}}

\newcommand{\phasecard}[4]{
\begin{tcolorbox}[enhanced,colback=#3,colframe=#2,arc=3mm,boxrule=.5pt,height=3.05cm,valign=center,halign=center,left=1.5mm,right=1.5mm,top=2mm,bottom=2mm]
{\sffamily\bfseries\fontsize{22}{24}\selectfont\textcolor{#2}{#1}}\par
{\sffamily\bfseries\small #4}
\end{tcolorbox}
}

% ── Recursos visuales de la guía ──────────────────────────────────────
% Declarados en el bloque `recursos:` del YAML y emitidos por el builder.
% \guiaFigura   → imagen rasterizada o PDF (foto, carta celeste, montaje).
% \guiaDiagrama → fuente TikZ/circuitikz (.tex) incrustada como vector:
%                 es la vía para circuitos y esquemáticos editables.
% #1 = ruta relativa al .tex · #2 = pie (puede ir vacío)
\newcommand{\guiaFigura}[2]{%
\begin{center}
\includegraphics[width=0.88\linewidth,height=0.38\textheight,keepaspectratio]{#1}\\[1.5mm]
{\footnotesize\itshape\textcolor{milcNegro!75}{#2}}
\end{center}
}

\newcommand{\guiaDiagrama}[2]{%
\begin{center}
\input{#1}\\[1.5mm]
{\footnotesize\itshape\textcolor{milcNegro!75}{#2}}
\end{center}
}

\begin{document}
\thispagestyle{empty}
\begin{tikzpicture}[remember picture,overlay]
  \fill[white] (current page.south west) rectangle (current page.north east);
  \path[fill=milcGradePrimary,rounded corners=16mm]
    ([xshift=.55cm,yshift=-.55cm]current page.north west) rectangle
    ([xshift=-.55cm,yshift=3.65cm]current page.south east);

  \node[anchor=north west,text=milcGradeText] at ([xshift=2.0cm,yshift=-1.85cm]current page.north west)
    {\sffamily\bfseries\fontsize{14}{16}\selectfont MOMENTO};
  \node[anchor=north west,text=milcGradeText,opacity=.82] at ([xshift=2.0cm,yshift=-2.75cm]current page.north west)
    {\sffamily\bfseries\fontsize{9}{11}\selectfont BEBRAS · PENSAR COMO UNA MÁQUINA};

  \begin{scope}[shift={([xshift=-3.05cm,yshift=-2.55cm]current page.north east)}]
    \fill[white,opacity=.80] (-.62,-.52) rectangle (.62,.52);
    \draw[milcGradeText,line width=1pt] (-.62,-.52) rectangle (.62,.52);
    \fill[milcGradeDark] (-.42,-.38) rectangle (-.22,.06);
    \fill[milcGradeAccent] (-.10,-.38) rectangle (.10,.24);
    \fill[milcGradePrimary!60!black] (.22,-.38) rectangle (.42,.38);
  \end{scope}

  \node[anchor=west,text=milcGradeText] at ([xshift=1.9cm,yshift=-12.85cm]current page.north west)
    {\sffamily\bfseries\fontsize{124}{124}\selectfont 1};
  % El circulito de grado (°) solo aplica a las guias de grado («11°»); en
  % Bebras y Semillero el numero grande es un momento/modulo, no un grado.
  

  \node[anchor=west,text=milcGradeText,text width=8.7cm,align=left] at ([xshift=10.35cm,yshift=-11.6cm]current page.north west)
    {
     {\sffamily\bfseries\fontsize{19}{22}\selectfont Parte y vencerás --- la descomposición\\como primera pieza del pensamiento computacional}\\[4mm]
     {\sffamily\bfseries\fontsize{23}{26}\selectfont Bebras · Momento 1}\\[2mm]
     {\sffamily\fontsize{13}{16}\selectfont Curso «Pensar como una máquina» · Pensamiento computacional}\\[5mm]
     {\sffamily\bfseries\fontsize{11}{13}\selectfont Institución Educativa Sor María Juliana}\\[1mm]
     {\sffamily\fontsize{10}{12}\selectfont Autor: Dr. Álvaro Cárdenas Orozco}
    };

  \path[fill=milcGradeSoft,rounded corners=7mm]
    ([xshift=1.35cm,yshift=.85cm]current page.south west) rectangle
    ([xshift=-1.35cm,yshift=4.25cm]current page.south east);
  \draw[milcGradeDark,line width=.65pt,rounded corners=7mm]
    ([xshift=1.35cm,yshift=.85cm]current page.south west) rectangle
    ([xshift=-1.35cm,yshift=4.25cm]current page.south east);
  \node[anchor=center,text=milcNegro,text width=17.0cm,align=center] at ([yshift=2.55cm]current page.south)
    {\sffamily\fontsize{10}{12}\selectfont Metodología MILC · Apertura ancestral · Pensamiento computacional · Triángulo Dussel-Estoicismo-Floridi\\
    \textcolor{milcGradeDark}{\bfseries Producto:} Tu problema real partido en 4 a 5 partes con su nivel de dificultad, más el plan de pasos de una meta tuya convertida en checklist.};
\end{tikzpicture}
\mbox{}
\newpage

\section*{Apertura: Parte y vencerás: la descomposición como primera pieza del pensamiento computacional}

\begin{softbox}{milcGradeDark}{milcGradeSoft}{Saber ancestral}
En el Litoral Pacífico, las artesanas del pueblo Wounaan tejen canastos en werregue tan apretados que retienen agua sin gotear. Esa obra casi imposible no nace de un solo tirón: la maestra la parte en tareas pequeñas y completas. Primero corta y tiñe la fibra de la palma de werregue. Luego arma el fondo en espiral, vuelta sobre vuelta, sin afán. Después levanta la pared, una hilera a la vez, sin saltarse ninguna. Y al final cierra el borde para que nada se deshaga. Cada vuelta es un trabajo terminado en sí mismo; nadie teje el canasto entero de golpe. Lo que de lejos parece magia, de cerca es \textbf{paciencia ordenada}: una obra enorme se vuelve posible cuando se parte en partes que sí caben en las manos.
\end{softbox}

\begin{softbox}{milcTurquesa}{milcGris}{Saber contemporáneo}
\textbf{Descomponer} es la primera y más poderosa de las cuatro piezas del pensamiento computacional: partir un problema grande en partes pequeñas que sí puedes resolver, atacarlas una a una y luego combinarlas. ¿Por qué funciona? Porque \textbf{lo enorme paraliza, pero lo pequeño se ataca}. Un computador nunca resuelve un problema gigante de golpe: lo corta en pasos y los hace en orden. En Bebras, muchas tareas parecen imposibles a primera vista\ldots hasta que las partes. Descomponer tiene \textbf{tres movimientos}: (1) identifica las partes, (2) resuelve cada una por separado, (3) combina las soluciones y revisa que encajen. Cuando las partes son independientes y se hacen en paralelo, los tiempos no se suman: terminas cuando acaba la parte más larga.
\end{softbox}

\begin{softbox}{milcMostaza}{milcGris}{Pregunta-puente}
\textit{La maestra wounaan no teje el canasto completo de un golpe: lo parte en vueltas pequeñas que sí puede tejer, una a la vez. ¿Y si te dijera que un computador --- y tú frente a un examen difícil --- resuelven sus retos exactamente así, partiéndolos en pedazos manejables?}
\end{softbox}

\begin{softbox}{milcVerde}{milcGris}{Saber y saber hacer}
Al terminar podrás: (1) \textbf{identificar} las partes pequeñas que componen un problema grande; (2) \textbf{explicar} por qué partir un problema lo vuelve más fácil de resolver; (3) \textbf{analizar} un problema tuyo descomponiéndolo y hallando su parte más difícil; (4) \textbf{aplicar} la descomposición para convertir una meta en un plan de pasos.
\end{softbox}



\small
\begin{tabularx}{\textwidth}{>{\bfseries}p{3.0cm}Y}
\rowcolor{milcGradePrimary!12}Autor & Dr. Álvaro Cárdenas Orozco\\
DBA & Desarrolla el pensamiento computacional --- descomposición, patrones, abstracción y algoritmos --- para resolver problemas (transversal 6°--11°, alineado a los lineamientos MEN de Tecnología e Informática)\\
Referentes & Valentina Dagienė (Bebras) · Pensamiento computacional (Wing) · Dussel · Estoicismo · Floridi · Saberes del Valle y el Pacífico\\
Producto final & Tu problema real partido en 4 a 5 partes con su nivel de dificultad, más el plan de pasos de una meta tuya convertida en checklist.\\
\end{tabularx}
\normalsize

\puente{Ya sabes que la artesana wounaan parte el canasto en vueltas posibles. Ahora mira el mapa de la sesión: vas a recorrer esos mismos tres movimientos de la descomposición hasta resolver tu primer reto de este tipo.}

\begin{titlebox}{milcGradeDark}
Ruta MILC: así vamos a aprender
\end{titlebox}
\begin{tabularx}{\textwidth}{>{\centering\arraybackslash}X>{\centering\arraybackslash}X>{\centering\arraybackslash}X>{\centering\arraybackslash}X}
\phasecard{1}{milcVerde}{milcGris}{Escucha\\[-1mm]\small Observo, escucho y pregunto.} &
\phasecard{2}{milcTurquesa}{milcGris}{Sistematización\\[-1mm]\small Pensamiento computacional.} &
\phasecard{3}{milcMagenta}{milcGris}{Praxis\\[-1mm]\small Construyo y aplico.} &
\phasecard{4}{milcVino}{milcGris}{Evaluación\\[-1mm]\small Reflexión liberadora.}
\end{tabularx}


\puente{Antes de nombrar la pieza, conviene verla en acción. Observa una tarea grande de tu entorno: la descomposición ya está ahí, aunque nadie le ponga ese nombre.}

\begin{titlebox}{milcVerde}
Fase 1 · Escucha
\end{titlebox}

\begin{softbox}{milcVerde}{milcGris}{Escena para escuchar}
\textbf{Observa una tarea grande} de tu casa o del colegio --- organizar el bazar, pintar el mural del patio, preparar un cumpleaños, limpiar y reorganizar tu cuarto. Fíjate cómo esa tarea, que al verla completa parece imposible, se parte en pedazos: quién hace qué, qué se puede hacer al mismo tiempo, qué hay que terminar antes de empezar lo siguiente. Sin saberlo, quien organiza eso ya está descomponiendo, igual que la artesana wounaan teje vuelta por vuelta.
\end{softbox}

\checkbox Nombraste la tarea grande y la partiste en al menos cuatro partes pequeñas y concretas.\\
\checkbox Señalaste cuáles partes pueden hacerse al mismo tiempo (en paralelo) y cuáles deben ir en orden.\\
\checkbox Identificaste qué parte hay que terminar primero para que las demás puedan avanzar.

\begin{infoband}{milcVerde}


\textbf{En tu cuaderno (Actividad 1 · IDENTIFICA).} Título: «Actividad 1 --- La tarea que descompongo». \textbf{Formato:} lista de 4 a 6 viñetas. \textbf{Extensión:} una viñeta por parte identificada, con el nombre de la parte y quién la haría. \textbf{Sabes que terminaste cuando} tu lista nombra, en lenguaje cotidiano, al menos cuatro partes distintas de la tarea grande que observaste.
\end{infoband}

\puente{Acabas de ver la descomposición en lo cotidiano. Ahora vamos a nombrar sus movimientos y entender por qué funcionan, porque con ellos vas a resolver todo lo que sigue.}

\begin{titlebox}{milcTurquesa}
Fase 2 · Sistematización con pensamiento computacional
\end{titlebox}

\begin{softbox}{milcTurquesa}{milcGris}{Cuatro pilares aplicados al tema}
Descomponer no es magia ni requiere código: es un gesto que se repite siempre igual, como la artesana teje vuelta por vuelta. Tiene \textbf{tres movimientos} que se aplican a cualquier problema. Pero la descomposición no trabaja sola: forma parte de las cuatro piezas del pensamiento computacional. Aquí las ves todas, con la descomposición en el centro porque es el punto de partida.
\end{softbox}

\small
\begin{tabularx}{\textwidth}{>{\bfseries\RaggedRight\arraybackslash}p{3.6cm}Y}
\rowcolor{milcTurquesa!90}\color{white}Pilar & \color{white}\bfseries Aplicación al tema\\
1. Descomposición & \textbf{Descomposición (esta pieza):} parte el problema grande en partes pequeñas, resuelve cada una por separado y combina las soluciones. Organizar el bazar del colegio se vuelve: decoración, comida, música y recaudo --- cuatro partes que sí caben en las manos.\\
2. Reconocimiento de patrones & \textbf{Reconocer patrones:} busca lo que se repite para no resolver lo mismo dos veces. Si cada mañana tardas el mismo tiempo en alistar el morral, ya puedes predecir cuándo salir sin mirar el reloj.\\
3. Abstracción & \textbf{Abstracción:} quédate con los datos que cambian la respuesta e ignora el resto. Cuando partes el bazar en áreas, no necesitas saber el color exacto de cada globo: solo cuántos equipos hay y qué hace cada uno.\\
4. Algoritmo conceptual & \textbf{Algoritmo:} escribe los pasos exactos en el orden correcto. Una vez que tienes las partes identificadas, el algoritmo dice en qué orden atacarlas: primero lo que desbloquea lo demás, luego lo que puede ir en paralelo, al final el ensamble.\\
\end{tabularx}
\normalsize

\begin{drawbox}{milcTurquesa}{Los 3 movimientos de la descomposición}
\textbf{1. Identifica las partes:} mira el problema grande y pregúntate ¿de qué pedazos está hecho? Como la artesana separa fibra, fondo, pared y borde. \\[1mm] \textbf{2. Resuelve cada parte:} ataca una parte a la vez, como una sola vuelta del canasto. Cada parte es pequeña, así que sí cabe en las manos: el «no puedo con todo» se vuelve «puedo con esta». \\[1mm] \textbf{3. Combina las soluciones:} junta las partes resueltas y revisa que encajen. Como cuando los equipos del bazar terminan lo suyo y el evento entero queda armado.
\end{drawbox}

\begin{softbox}{milcMostaza}{milcGris}{Errores comunes y su corrección}
\textbf{Querer resolver todo de un golpe:} el problema parece imposible solo porque no lo has partido. \textbf{Confundir las partes con los pasos:} las partes son piezas del problema; los pasos son el orden para resolverlas. \textbf{Olvidar combinar al final:} resolver las partes sin ensamblarlas deja el trabajo incompleto, como tejer el fondo y la pared por separado sin unirlos nunca.
\end{softbox}

\begin{infoband}{milcTurquesa}


\textbf{En tu cuaderno (Actividad 2 · ANALIZA).} Título: «Actividad 2 --- Los 3 movimientos en mi tarea». \textbf{Formato:} tabla de 2 columnas (movimiento · cómo lo apliqué yo). \textbf{Extensión:} los tres movimientos aplicados a la misma tarea que observaste en la Actividad 1. \textbf{Sabes que terminaste cuando} cada celda tiene un ejemplo concreto tuyo, no copiado del texto.
\end{infoband}

\puente{Ya distingues los tres movimientos de la descomposición. Es hora de usarlos: vas a resolver un reto estilo Bebras aplicándolos a propósito, paso a paso.}

\begin{titlebox}{milcMagenta}
Fase 3 · Praxis con pensamiento computacional
\end{titlebox}

\begin{softbox}{milcMagenta}{milcGris}{Construyendo el producto con los cuatro pilares}
Los movimientos no son adorno: se usan. Vas a resolver un \textbf{reto estilo Bebras sobre descomposición} aplicando los tres movimientos a propósito. Lee con calma y parte el problema antes de responder --- la clave está en descomponer, no en adivinar.
\end{softbox}

\small
\begin{tabularx}{\textwidth}{>{\bfseries\RaggedRight\arraybackslash}p{3.6cm}Y}
\rowcolor{milcMagenta!90}\color{white}Pilar & \color{white}\bfseries Aplicación al producto\\
Descomposición & \textbf{Descomposición:} parte el enunciado. ¿Cuántas partes hay? ¿Qué hace cada una? ¿Son independientes o una depende de otra?\\
Patrones & \textbf{Patrones:} busca si algo se repite --- una regla de conteo, una estructura que se repita en filas o columnas.\\
Abstracción & \textbf{Abstracción:} ignora los datos de adorno; quédate con los números y relaciones que cambian la respuesta.\\
Algoritmo de armado & \textbf{Algoritmo:} escribe, en orden, los pasos que te llevan del enunciado a la respuesta, sin saltarte ninguno.\\
\end{tabularx}
\normalsize

\begin{drawbox}{milcMagenta}{El reto: el mural en cuadrícula}
\checkbox \textbf{Reto (Nivel B).} El mural del patio se divide en una cuadrícula de \textbf{4 columnas por 3 filas}. Cada estudiante pinta exactamente \textbf{un cuadro}. ¿Cuántos estudiantes se necesitan para cubrir \emph{todos} los cuadros? \\[2mm]
\checkbox Marca: \quad \fbox{\phantom{x}}~7 \quad \fbox{\phantom{x}}~12 \quad \fbox{\phantom{x}}~4 \quad \fbox{\phantom{x}}~34 \\[2mm]
\checkbox Escribe \textbf{qué movimiento de la descomposición} usaste para resolverlo y por qué.
\end{drawbox}

\begin{softbox}{milcMagenta}{milcGris}{Plantilla de trabajo}
\textbf{Resuélvelo paso a paso (descomposición).} Parte el problema: la cuadrícula tiene \textbf{4 columnas} y \textbf{3 filas}. Cada columna tiene 3 cuadros; hay 4 columnas: $4 \times 3 = 12$ cuadros en total. Como cada estudiante pinta exactamente uno, se necesitan \textbf{12 estudiantes}. Movimiento usado: \textbf{identificar las partes} --- separaste filas y columnas, resolviste cuántos cuadros hay en total y obtuviste la respuesta.
\end{softbox}

\begin{infoband}{milcMagenta}


\textbf{En tu cuaderno (Actividad 3 · APLICA).} Título: «Actividad 3 --- Mi reto Bebras de descomposición». \textbf{Formato:} respuesta marcada + 3 renglones de explicación. \textbf{Extensión:} la respuesta y el razonamiento paso a paso, nombrando el movimiento usado. \textbf{Sabes que terminaste cuando} tu explicación convencería a alguien que aún no sabe la respuesta.
\end{infoband}

\puente{Resolviste el reto. Ahora deja tu evidencia: el problema partido en partes y el checklist de tu meta, para mostrar que sabes descomponer de verdad.}

\begin{titlebox}{milcMostaza}
Producto de la sesión
\end{titlebox}
\textbf{Problema real partido en partes con su nivel de dificultad, más checklist de meta personal}.
\begin{softbox}{milcMostaza}{milcGris}{Criterios de calidad}
(1) Partiste un problema real en 4 a 5 partes pequeñas y concretas, con el nivel de dificultad de cada una. (2) Explicaste cuál es la parte más difícil y por qué cuesta más que las otras. (3) Tu checklist de meta tiene 5 a 7 pasos en orden lógico, del primero al último. (4) Resolviste el reto del mural y marcaste la respuesta correcta (12 estudiantes). (5) Nombraste el movimiento de la descomposición que usaste para resolver el reto, con una frase que explica cómo lo aplicaste.
\end{softbox}

\puente{Terminaste el producto. Detente a mirar qué cambió en ti --- no la nota, sino tu manera de enfrentar lo que parecía imposible.}

\begin{titlebox}{milcVino}
Fase 4 · Evaluación liberadora — cinco dimensiones
\end{titlebox}

\begin{softbox}{milcVino}{milcGris}{Sentido de la evaluación}
Evaluar no es solo revisar una nota. Es reconocer qué aprendí, qué cambió en mí, qué emociones manejé, cómo me relaciono con la ciudad y la comunidad, y qué le dejaré a quien viene detrás.
\end{softbox}

\textbf{Mi reflexión liberadora en cinco dimensiones.} Completa cada frase iniciada:

\small
\begin{tabularx}{\textwidth}{>{\bfseries\RaggedRight\arraybackslash}p{3.4cm}Y>{\RaggedRight\arraybackslash}p{4.6cm}}
\rowcolor{milcVino}\color{white}Dimensión & \color{white}\bfseries Mi respuesta (frame starter) & \color{white}\bfseries Algo que descubrí\\
1. Personal & Lo que más cambió en mí fue \linead{6.5cm} & \linead{4.2cm}\\
2. Emocional & La emoción más fuerte fue \linead{2.5cm} y la regulé \linead{3cm} & \linead{4.2cm}\\
3. Ciudadana & En democracia esto importa porque \linead{6cm} & \linead{4.2cm}\\
4. Local (Cartago) & En mi barrio sería útil para \linead{6.5cm} & \linead{4.2cm}\\
5. Intergeneracional & A mi abuela/abuelo le diría \linead{6.5cm} & \linead{4.2cm}\\
\end{tabularx}
\normalsize

\begin{infoband}{milcVino}
\textbf{Para la próxima sustentación.} Vuelve a esta tabla en seis meses. Verás que las respuestas cambian — no porque lo aprendido cambie, sino porque tú cambias. La evaluación liberadora es una conversación contigo a través del tiempo.
\end{infoband}


\puente{Cerramos pensando en grande: tres voces para entender qué significa, para ti y para tu comunidad, aprender a partir lo enorme en pedazos posibles.}

\begin{titlebox}{milcGradeDark}
Triángulo de pensamiento · cierre filosófico
\end{titlebox}

\begin{softbox}{milcVino}{milcGris}{Dussel · lente del nosotros}
\textit{``El saber del pueblo también es ciencia.''}

El pueblo Wounaan ya descomponía obras enormes en partes --- fondo, pared, borde --- mucho antes de que existiera la palabra «algoritmo». La descomposición no llegó de afuera: tu comunidad la practicaba tejiendo, construyendo y organizando mingas desde siempre.

\textbf{¿Qué saberes de mi familia o mi comunidad ya partían tareas imposibles en pedazos posibles, aunque nadie los llamara pensamiento computacional?}
\end{softbox}
\linead{\linewidth}\\[1mm]
\linead{\linewidth}

\begin{softbox}{milcMostaza}{milcGris}{Estoicismo · Marco Aurelio · cuidado interior}
\textit{``Haz cada acto de tu vida como si fuera el último.''}

Lo enorme asusta y paraliza. Pero si te concentras en hacer bien \emph{una sola parte} a la vez --- con atención plena, como la artesana en su vuelta --- el «no puedo con todo esto» se transforma en «puedo con esta parte ahora».

\textbf{¿Puedo elegir una sola parte del problema que me abruma y hacerla con toda mi atención, en vez de bloquearme mirando el todo?}
\end{softbox}
\linead{\linewidth}\\[1mm]
\linead{\linewidth}

\begin{softbox}{milcTurquesa}{milcGris}{Floridi · lente de la infoesfera}
\textit{``Somos inforgs: organismos informacionales que habitan la infoesfera.''}

Un computador nunca procesa un problema gigante de golpe: lo descompone en instrucciones pequeñas y las ejecuta en orden. Aprender a descomponer es entender cómo «piensa» la máquina con la que ya convives --- y descubrir que tú también lo haces, sin darte cuenta.

\textbf{Si un computador resuelve lo difícil partiéndolo en pasos pequeños, ¿en qué se parece mi forma de pensar a la suya cuando descompongo un problema mío?}
\end{softbox}
\linead{\linewidth}\\[1mm]
\linead{\linewidth}

\begin{titlebox}{milcVino}
Compromiso final
\end{titlebox}

\small
\begin{tabularx}{\textwidth}{>{\RaggedRight\arraybackslash}p{3.0cm}YYY}
\rowcolor{milcVino}\color{white}\bfseries Criterio & \color{white}\bfseries Lo logré & \color{white}\bfseries En proceso & \color{white}\bfseries Necesito apoyo\\
Comprensión del tema & Explico ideas centrales con mis palabras. & Reconozco ideas pero falta precisión. & Necesito repasar conceptos base.\\
Pensamiento computacional & Apliqué los 4 pilares al diseñar mi producto. & Apliqué algunos pilares. & Necesito releer la fase 2.\\
Aplicación y producto & Mi producto cumple los criterios y se puede revisar. & Producto funcional con ajustes pendientes. & Aún en construcción.\\
Reflexión 5 dimensiones & Respondí cada dimensión con honestidad. & Respondí algunas con detalle. & Mis respuestas son muy generales.\\
Triángulo de pensamiento & Conecté las 3 voces con mi trabajo real. & Conecté al menos una. & No relaciono las voces con mi proceso.\\
\end{tabularx}
\normalsize

\textbf{Hoy entendí que:}\\[1mm]
\linead{\linewidth}\\[1mm]
\linead{\linewidth}

\textbf{Mi compromiso para la próxima sesión es:}\\[1mm]
\linead{\linewidth}\\[1mm]
\linead{\linewidth}

\begin{softbox}{milcMostaza}{milcGris}{Cita de despedida · Marco Aurelio}
\textit{``El obstáculo en el camino se vuelve el camino. Lo que estorba al camino se vuelve camino.''}\\[1mm]
Cada error de hoy, bien observado, se convierte en pieza del camino que pisarás mañana.
\end{softbox}

\small
\begin{tabularx}{\textwidth}{>{\bfseries}p{3.6cm}Y}
\rowcolor{milcGradeDark!12}Nombre completo: & \linead{10cm}\\
Fecha: & \linead{4cm} \quad Curso/grupo: \linead{4cm}\\
Mi firma: & \linead{9cm}\\
\end{tabularx}
\normalsize

\begin{infoband}{milcGradeDark}
\textbf{Para la próxima sesión.} Elige una tarea que te parezca imposible esta semana --- un examen que se viene, un proyecto pendiente, un cuarto que no has ordenado --- y pártela en partes pequeñas en tu cuaderno. Escribe el nombre de cada parte, cuál harías primero y por qué. Trae tu lista para compartirla con el grupo: vamos a ver cuántas partes distintas encontró cada uno en el mismo problema.
\end{infoband}

\end{document}
